Large Language Models (LLMs) have significantly advanced natural language processing, but content repetition in open-source LLMs remains a critical challenge that adversely affects user experience. The repetition penalty parameter (\rpp) aims to mitigate this issue by preventing repeated content generation, but excessive use of \rpp can compromise the overall quality.
In this paper, we propose Repetition-Aware Performance (\rap), a novel evaluation metric that quantifies and integrates repetition penalty into the assessment of model performance, enabling tuning of \rpp. 
We evaluate our approach using twelve open-source LLMs, ranging from 2 billion to 70 billion parameters, tested on question answering and machine translation tasks across three datasets with varying prompting techniques. 
Experimental results show that RAP effectively tunes RPP, helping to identify a trade-off value that significantly reduces repetition while minimizing performance loss.
The code and the dataset of generated text can be accessed at \url{https://github.com/inflaton/rap}.
% Upon acceptance, we will release the code and the dataset of generated text, providing a valuable resource for further research on repetition detection and LLMs evaluation.

% Large Language Models (LLMs) have significantly advanced the field of natural language processing; however, the issue of content repetition in open-source LLMs remains a critical challenge, negatively affecting user experience. The repetition penalty parameter (\rpp) is designed to mitigate this issue by preventing the generation of repeated content, but excessive application of \rpp can compromise the overall quality of the generated text. 
%

% Additionally, we find that chat templates effectively reduce repetition, a strategy underused in LLM applications development. 
% We assess our approach using twelve LLMs, with sizes ranging from 2 billion to 70 billion parameters, tested on question answering and machine translation tasks, utilizing three datasets and varying prompting techniques.
% These techniques include both Retrieval-Augmented Generation (RAG) and non-RAG strategies, using generic prompt and chat template. 
% We evaluate twelve LLMs, with sizes ranging from 2 billion to 70 billion parameters, yielding insights into optimal parameter settings for minimizing repetition while maintaining quality.
%
% Experimental results show that increasing \rpp generally decreases performance, thus tuning \rpp to find the balance is important. 
% Our findings show that RAP can be used to effectively tune RPP, helping to identify a trade-off value that significantly reduces repetition while minimizing performance loss.
% We also find that chat templates effectively reduce repetition, a strategy underused in LLM applications development. 
% Upon acceptance, we will release the code and the dataset of generated text, providing a valuable resource for further research on repetition detection and LLMs evaluation.

% We also find that increasing \rpp generally reducing performance, therefore, having a balance approach is desireable, that can be tuned using our proposed metric, \rap. 
% Upon acceptance, we will make the source code and all experimental data publicly available, including a substantial corpus of text generated by open-source LLMs. This dataset will serve as a valuable resource for future research in developing repetition detection methods and evaluating LLM capabilities in question-answering tasks.


% Large Language Models (LLMs) have significantly advanced natural language processing, yet content repetition issue in open-source LLMs remains a challenge and negatively impacts user experience in applications like conversational chatbots. Repetition penalty parameter (\rpp) is designed to stop LLMs from generating repeated content. However, excessive use of \rpp can impact the quality of the generated text. In this paper, we propose 

% \todo{to revise, add findings, add MT}
% We propose Repetition-Aware Performance (\rap) metric, to investigate the impact of repetition penalty parameter (\rp) on the quality of responses generated by open-source LLMs. RAPGeT incorporates a novel metric, the repetition factor (RF), to quantify the repetition penalty into model's performance. This enable identifying the trade-off between repetition and overall quality.
% Our evaluation on two question-answering datasets (\webqsp and \ms), using both Retrieval-Augmented Generation (RAG) and non-RAG approaches, provides insights into optimal parameter settings. 
% We analyze nine LLMs with sizes ranging from 2 billion to 70 billion parameters, providing insights on reducing repetition while maintaining quality in open-source LLMs.
% Our findings show that an optimal balance between repetition suppression and response quality can be achieved with careful tuning of the RPP.
% Upon acceptance, we will publish the source code and our experiment results. This data can be used for future tasks such as designing repetition detection methods or evaluating LLMs’ capabilities in QA tasks.


% Large Language Models (LLMs) have significantly advanced natural language processing, yet content repetition issue in open-source LLMs remains a challenge and negatively impacts user experience in applications like conversational chatbots. The \textit{repetition penalty parameter} (\rp) mitigates this issue by penalizing frequent tokens but can lead to incomplete responses if applied stringently.
% %
% In this paper, we propose Repetition-Aware Performance Evaluation for Generated Text (\ours) framework, to investigate \rp effect on the quality of responses generated by open-source LLMs.
% We define the problem of repetition penalty detection and introduce the \textit{repetition factor}, a novel metric that quantifies repetition penalty on performance, allowing us to identify the trade-off between repetition and overall quality.
% Our evaluation on two question-answering datasets (\webqsp and \ms), using both Retrieval-Augmented Generation (RAG) and non-RAG approaches, provides insights into optimal parameter settings. 
% We analyze nine LLMs with sizes ranging from 2 billion to 70 billion parameters, providing insights on reducing repetition while maintaining quality in open-source LLMs.
% Our findings show that an optimal balance between repetition suppression and response quality can be achieved with careful tuning of the RPP. For instance, models like Llama-2-7B and Llama-3-8B experience significant improvements in repetition-aware performance scores when using the RAPGeT framework, demonstrating the effectiveness of our approach. 
% Upon acceptance, we will publish the source code and our experiment results. This data can be used for future tasks such as designing repetition detection methods or evaluating LLMs’ capabilities in QA tasks. 

% Large Language Models (LLMs) have significantly advanced natural language processing, yet content repetition issue in open-source LLMs remains a challenge and negatively impacts user experience in applications
% We propose Repetition-Aware Performance Evaluation for Generated Text (\ours) framework, to investigate \rp effect on the quality of responses generated by open-source LLMs.
% We define the problem of repetition penalty detection and introduce the \textit{repetition factor}, a novel metric that quantifies repetition penalty on performance, allowing us to identify the trade-off between repetition and overall quality.
% Our evaluation on two question-answering datasets (\webqsp and \ms), using both Retrieval-Augmented Generation (RAG) and non-RAG approaches, provides insights into optimal parameter settings. 
% We analyze nine LLMs with sizes ranging from 2 billion to 70 billion parameters, providing insights on reducing repetition while maintaining quality in open-source LLMs.
% Our findings show that an optimal balance between repetition suppression and response quality can be achieved with careful tuning of the RPP. For instance, models like Llama-2-7B and Llama-3-8B experience significant improvements in repetition-aware performance scores when using the RAPGeT framework, demonstrating the effectiveness of our approach. 
% Upon acceptance, we will publish the source code and our experiment results. This data can be used for future tasks such as designing repetition detection methods or evaluating LLMs’ capabilities in QA tasks. 


% This paper introduces the Repetition-Aware Performance Evaluation for Generated Text (RAPGeT) framework to address content repetition in open-source Large Language Models (LLMs), particularly impacting user experience in conversational applications. RAPGeT incorporates a novel metric, the repetition factor (RF), to quantify the repetition penalty into model's performance and mitigate repetition through the adjustment of the repetition penalty parameter (RPP). Our comprehensive evaluation across two major question-answering datasets—WebQSP and MS MARCO—utilizing both Retrieval-Augmented Generation (RAG) and non-RAG approaches, reveals that fine-tuning RPP effectively balances repetition suppression with response quality. This framework demonstrates significant performance improvements in models like Llama-2-7B and Llama-3-70B, illustrating its potential to enhance the quality of responses in LLMs without sacrificing content richness.



% We propose Repetition-Aware Performance Evaluation for Generated Text (\ours) framework, to investigate the impact of repetition penalty parameter (\rp) on the quality of responses generated by open-source LLMs. RAPGeT incorporates a novel metric, the repetition factor (RF), to quantify the repetition penalty into model's performance. This enable identifying the trade-off between repetition and overall quality.
% Our evaluation on two question-answering datasets (\webqsp and \ms), using both Retrieval-Augmented Generation (RAG) and non-RAG approaches, provides insights into optimal parameter settings. 
% We analyze nine LLMs with sizes ranging from 2 billion to 70 billion parameters, providing insights on reducing repetition while maintaining quality in open-source LLMs.
% Our findings show that an optimal balance between repetition suppression and response quality can be achieved with careful tuning of the RPP.
% Upon acceptance, we will publish the source code and our experiment results. This data can be used for future tasks such as designing repetition detection methods or evaluating LLMs’ capabilities in QA tasks. 

% Our findings show that an optimal balance between repetition suppression and response quality can be achieved with careful tuning of the RPP. For instance, models like Llama-2-7B and Llama-3-8B experience significant improvements in repetition-aware performance scores when using the RAPGeT framework, demonstrating the effectiveness of our approach. 

% Our findings show that an optimal balance between repetition suppression and response quality can be achieved with careful tuning of the RPP. For instance, models like Llama-2-7B and Llama-3-8B experience significant improvements in repetition-aware performance scores when using the RAPGeT framework, demonstrating the effectiveness of our approach. 
% Upon acceptance, we will publish the source code and our experiment results. This data can be used for future tasks such as designing repetition detection methods or evaluating LLMs’ capabilities in QA tasks. 